\dykband{EXERCISES IN SWAFL}
{\small Open SWAFL (desktop or online) with the benchmark project. The chapter of the \href{https://github.com/halilunaziru73-creator/Coupled-Water-Fluxes-Pore-Scale-Dynamics-Precision-Irrigation-Mulching-Feedback-Loops/raw/main/Software/v1.1.0/SWAFL_Desktop/docs/manual_v1.1.0/SWAFL_Manual_v1.1.0.pdf}{SWAFL manual} that explains each task is given in brackets; record your values in a table and compare them with what is expected (\emph{Check}).}\par\vspace{3pt}
\begin{exlist}
\item {\color{khblue}[Manual, Chapter 8]} Open the layer \emph{Soil structure and tillage} and choose the sand and then the silty clay loam. For each, read the saturated water content and sketch the retention curve that the layer draws. Which soil has the wider range of pore sizes?\par\noindent{\small\itshape Check: the silty clay loam (smaller $n$), whose curve falls gradually over several decades of suction.}
\item {\color{khblue}[Manual, Chapter 8]} In the same layer move the \emph{Structural consolidation rate} to its smallest and largest value and run. Record how the drainage and the evaporation of the year change, and explain the change with the pore picture of this chapter.\par\noindent{\small\itshape Check: faster consolidation removes the large pores made by tillage, so the drainage of the loosened layer falls.}
\item {\color{khblue}[Manual, Chapter 9]} Open \emph{Mulch and evaporation} and watch the pore picture and the drying front during \emph{Play}. Note the day on which stage 2 begins on the bare soil, with and without mulch.\par\noindent{\small\itshape Check: stage 2 begins later under the mulch; the drying front stays near the surface.}
\item {\color{khblue}[Manual, Chapter 7]} In \emph{Soil, plant and air} compare the water potential of the soil, the root, the leaf and the air on a hot day. Which step has the largest drop, and why?\par\noindent{\small\itshape Check: the step from the leaf to the air, which is far larger than all the others.}
\end{exlist}
